%%%%%%%%%%%%%%%%%%%%%%%%%%%%%%%%%%%%%%%%%%%%%%%%%%%
%
%  CHAPTER VIII -- PROPOSED WORK II: MECHANICAL EXPLORATION
%
%%%%%%%%%%%%%%%%%%%%%%%%%%%%%%%%%%%%%%%%%%%%%%%%%%%

\chapter{\uppercase{Proposed Work II: Mechanical Exploration of cP-Weave Plate Lattices}}
\label{cha:mechanical}

The premise of this dissertation is that the arrangement of parts, not the parts themselves, determines the behavior of a fibrillar structure. This chapter proposes a direct physical test of that premise. We print cP-weaves as plate lattices with the tubular negative-space method of Chapter~\ref{cha:manufacturing}, so that the solid walls between thread channels carry the load, and compare them in compression with triply periodic minimal surface (TPMS) lattices printed under the same conditions. The design framework defines the specimens, the enumeration tells us which are distinct, and the manufacturing method makes them reproducible.

\section{A Deliberately Small Design Space}
\label{sec:mech-designspace}

The cP design space grows too quickly to test exhaustively (Chapter~\ref{cha:volumetric-bravais}). We therefore restrict it to a set in which only the arrangement varies. We use the 64 plain $2\times2\times2$ translation cells, each tiled to a $4\times4\times4$ domain of $4\,\mathrm{mm}$ cells, giving $16\,\mathrm{mm}$ cubes. In every plain cP-weave each strand has the same geometry and therefore the same volume within its cell, so one strand radius gives the same volume fraction for all 64 specimens. The specimens differ only in how the local configurations are arranged. Each is identified by a three-digit code, one digit per axis, from the compact representation of Section~\ref{sec:cp-weaves}.

The 64 cells fall into four orbits under the lattice isometry group. Specimens in the same orbit represent the same infinite weave up to translation, rotation, or reflection, but are cut from it differently or loaded along a different axis. Comparing them separates the effect of the weave from the effect of its orientation and boundary.

\section{Volume-Fraction Control}
\label{sec:vf-control}

A specimen is a solid cube from which tubes of radius $r$ around every thread centerline are removed. Increasing $r$ lowers the volume fraction $\mathrm{vf}$. The largest admissible radius is set by the minimum distance between neighboring strands, which we compute by measuring, for every point on every strand, the distance to all other strands. For the plain $1\_1\_1\_1$ configuration with weave offset $\mathrm{wo}=1.4603\,\mathrm{mm}$ and weave depth $\mathrm{wd}=0.0952\,\mathrm{mm}$, printer resolution narrows the useful range to $r\in[0.01,0.65]\,\mathrm{mm}$, giving $\mathrm{vf}\in[0.664,0.9999]$. We target $\mathrm{vf}=0.8$ and find $r$ by bisection to a tolerance of $0.001$ in $\mathrm{vf}$ (Figure~\ref{fig:vf-control}).

\begin{figure}[htbp]
\centering
\includegraphics[width=6in]{ch08_fig01_vf_control.png}
\caption[Volume-fraction control by strand radius.]{Volume-fraction control. (a) Weave offset, weave depth, and strand radius on one cell, with the minimum inter-strand clearance that bounds the radius. (b) Volume fraction as a function of strand radius, with the printable range and the target $\mathrm{vf}=0.8$.
\figtodo{No existing source. (a) New annotated render of one $4\,$mm plain $1\_1\_1\_1$ cell from the weave generator, marking wo, wd, $r$, and the closest-point segments to neighboring strands. (b) Plot of $\mathrm{vf}(r)$ from the bisection data in the printed-weave characterization notes: mark $r{=}0.65\,$mm$\to$0.664, $r{=}0.01\,$mm$\to$0.9999, and the $\mathrm{vf}{=}0.8$ target.}}
\label{fig:vf-control}
\end{figure}

As a preliminary check, we compared the target with the slicer-reported mass of four configurations. A solid $16\,\mathrm{mm}$ cube needs about $5.35\,\mathrm{g}$ of PLA, so $\mathrm{vf}=0.8$ predicts about $4.28\,\mathrm{g}$ (Table~\ref{tab:print-mass}). The mass-implied volume fractions are about one to two points above the target, which we attribute to the nozzle widening printed features. More importantly, the four configurations agree with each other to within $0.05\,\mathrm{g}$, or about $0.01$ in volume fraction, even though each radius was found independently. The method therefore holds volume fraction constant across configurations, which is the property the study needs.

\begin{table}[htbp]
\centering
\begin{tabular}{@{}lcc@{}}
\toprule
\textbf{Configuration} & \textbf{Slicer mass (g)} & \textbf{Mass-implied vf} \\
\midrule
$1\_1\_1\_1$ & 4.37 & 0.817 \\
$2\_1\_1\_1$ & 4.33 & 0.809 \\
$2\_2\_1\_1$ & 4.32 & 0.807 \\
$3\_1\_1\_1$ & 4.33 & 0.809 \\
\bottomrule
\end{tabular}
\caption{Slicer-reported mass and mass-implied volume fraction for four configurations at a target of $\mathrm{vf}=0.8$ (solid cube: $5.35\,$g).}
\label{tab:print-mass}
\end{table}

\section{Specimens and Testing}
\label{sec:tpms-specimens}

Two sets of $16\,\mathrm{mm}$ PLA cubes are printed at $\mathrm{vf}=0.8$ under the same conditions (Figure~\ref{fig:tpms-specimen-sets}):

\begin{itemize}
\item \textbf{Woven lattices}: the 64 plain cP-weave specimens of Section~\ref{sec:mech-designspace}.
\item \textbf{TPMS baselines}: eight families (gyroid, Schwarz P, Schwarz D, Schwarz F-RD, Schoen I-WP, Fischer-Koch S, Lidinoid, and split P), each at eight scales of its defining function. Changing the scale changes feature size at fixed volume fraction, so each family is represented by a range rather than one arbitrary specimen.
\end{itemize}

Print batch, build-plate position, and print and test dates are recorded and spread across both sets, so that a manufacturing effect cannot be mistaken for a design effect. Each cube is compressed uniaxially on a 10~kN frame with a stop at 9.5~kN. Raw exports are processed by an open analysis pipeline written for this study~\cite{companal2026}. It removes the pre-contact tare and seating slack, detects duplicate exports by content, and flags tests that reach the load limit, so that quantities depending on the unmeasured part of the curve are reported as lower bounds. Because displacement is crosshead travel, measured stiffness includes frame compliance, which we will measure separately for absolute moduli.

\begin{figure}[htbp]
\centering
\includegraphics[width=6in]{ch08_fig02_specimen_sets.png}
\caption[Woven and TPMS specimen sets.]{The two specimen sets: (top) woven lattices from the 64 plain cP-weave codes and (bottom) the eight TPMS families at one of their eight scales. All are $16\,\mathrm{mm}$ PLA cubes at $\mathrm{vf}=0.8$.
\figtodo{No existing source. Top: renders or photos of a few woven lattices, ideally one per isometry orbit. Bottom: one render per TPMS family at a common scale (names as in \fs{comp-anal/data/metadata/baseline_metadata.csv}).}}
\label{fig:tpms-specimen-sets}
\end{figure}

\section{Metrics and Analysis}
\label{sec:tpms-analysis}

From each stress-strain curve we extract standard cellular-solid metrics~\cite{gibson1997cellular,iso13314}: Young's modulus from a sliding-window fit, plateau stress, densification strain by the energy-efficiency criterion, absorbed and specific absorbed energy, first peak stress, and stress and energy at 10\% and 20\% strain. The fixed-strain metrics remain valid for every test that reaches those strains, including truncated ones.

The two sets are compared at three levels.
\begin{enumerate}
\item \textbf{Envelopes.} Curves are resampled on a common strain grid, and each class is summarized by its median and interquartile band, with the woven median shown in every TPMS panel (Figure~\ref{fig:tpms-envelopes}). The woven set is also split by each digit of the pattern code and by isometry orbit.
\item \textbf{Property space.} Pairs of metrics, such as modulus against plateau stress, show which regions each class reaches and whether the woven lattices occupy a region no TPMS family covers.
\item \textbf{Parameter-property statistics.} Spearman rank correlations with Benjamini-Hochberg control of the false discovery rate~\cite{benjamini1995fdr}, partial correlations to separate correlated parameters, and rank analysis of covariance for categorical factors such as family and batch~\cite{quade1967rank}. Gibson-Ashby exponents indicate whether a class is bending or stretching dominated.
\end{enumerate}
Finally, for each TPMS family we find the scale whose response is closest to the woven distribution, with a bootstrap interval, and compare it with the scales that are geometrically comparable to the woven lattices. Agreement would mean that a TPMS that looks like the woven lattices also behaves like them. Disagreement would point to an effect of the woven topology itself.

\begin{figure}[htbp]
\centering
\includegraphics[width=6in]{ch08_fig03_envelopes.png}
\caption[Planned stress-strain envelope comparison.]{Planned envelope comparison: one panel per TPMS family showing its median curve and interquartile band, with the woven-lattice median and band repeated in grey in every panel. The load limit is marked by a dashed line.
\figtodo{Planned figure from the \fs{comp-anal} pipeline~\cite{companal2026}, \fs{results/figures/class_envelopes.png} (\fs{scripts/run_analysis.py}). Until results are reported, use the synthetic demo (\fs{scripts/make_demo_data.py}) labeled ``illustrative'', or a schematic.}}
\label{fig:tpms-envelopes}
\end{figure}

\section{Expected Outcome}

The study will show whether rearranging the same strands at the same volume fraction moves the mechanical response as much as switching between TPMS families, and where woven lattices sit among established architected materials. A large, structured spread across the 64 arrangements, with consistent differences between orbits, would be direct evidence that in fibrillar structures the whole is greater than the sum of its parts. It also establishes a pipeline from a labeled design to a measured property that can be reused for other cell types and periodic H-GRS structures.
